
\begin{equation}\tag{3.1}
 c_i(f_*(y))=f_*\bigl(P_i((c_\alpha(y))_\alpha,(c_\beta(\underline N))_\beta)\bigr),
\end{equation}
où $P_i$ est un polynôme universel en les $c_\alpha(y)$, $c_\beta(\underline N)$, polynôme qui a priori est à coefficients rationnels. En fait, on constate [RRR] que ces polynômes sont à coefficients \uline{entiers}. Cela soulève la question de la validité d'une formule de Riemann--Roch ``sans dénominateurs''. En fait, une telle formule pour $c_i(f_*(y))$ ne peut être cherchée en général dans $\operatorname{Gr}^\bullet(X)$ lui-même, vu que l'homomorphisme $f_*:K^\bullet(Y)\to K^\bullet(X)$ n'est pas compatible lui-même avec les $\lambda$-filtrations modulo décalage par $d$ (4.6), mais l'est seulement à torsion près; de sorte qu'il n'est pas possible de définir un homomorphisme $f_*:\operatorname{Gr}^\bullet(Y)\longrightarrow \operatorname{Gr}^\bullet(X)$ donnant l'homomorphisme naturel $f_*:\operatorname{Gr}^\bullet(Y)_\Q\to\operatorname{Gr}^\bullet(X)_\Q$, et permettant de donner un sens à la formule (3.1) comme une formule dans $\operatorname{Gr}^\bullet(X)$, et non seulement dans $\operatorname{Gr}^\bullet(X)_\Q$. Mais lorsque $X$ est supposé régulier, donc $K^\bullet(X)\xrightarrow{\sim} K_\bullet(X)$ (IV 2.5), on peut munir $K^\bullet(X)$ de la \uline{filtration topologique} par la codimension des supports, héritée de la filtration correspondante de $K_\bullet(X)$. Il est plausible que cette filtration est compatible avec la structure d'anneau, fait qu'on peut prouver tout au moins [RRR] lorsque $X$ est lisse sur un corps $k$, en utilisant la théorie de l'anneau de Chow [4]. Quand il en est ainsi, le gradué associé
\[
 \Gtop^\bullet(X)
\]
est un \uline{anneau gradué}, et supposant également $Y$ régulier, on a fait ce qu'il fallait pour avoir un homomorphisme naturel de degré $d$ de groupes gradués
\[
 f_*=f_*^{\mathrm{top}}:\Gtop^\bullet(Y)\longrightarrow \Gtop^\bullet(X).
\]
Comme la $\lambda$-filtration de $K^\bullet(X)$ est plus fine que la filtration topologique par la codimension qu'on vient de considérer (X 1.3.3), on trouve un homomorphisme canonique de groupes gradués (et même d'anneaux gradués)
\[
 \operatorname{Gr}^\bullet(X)\longrightarrow \Gtop^\bullet(X),
\]
qui est d'ailleurs un isomorphisme mod torsion (Exp. VIII). Ceci permet, comme dans [RRR], de considérer une théorie des classes de Chern sur $X$ à valeurs dans $\Gtop^\bullet(X)$, et de même sur $Y$, et de donner ainsi un sens à la formule (3.1) comme une formule dans $\Gtop^\bullet(X)$. Sous cette forme, cette formule est prouvée dans [RRR] lorsque $X$ est de \uline{caractéristique nulle}, tout au moins si $X$ est lisse sur un corps (hypothèse servant uniquement à assurer que la filtration topologique de $K^\bullet(X)$ est compatible avec la structure d'anneau). L'hypothèse de caractéristique nulle est utilisée
essentielle%
